\clearpage
\fancyhead[L]{\small PASQUICH. EHRENRETTUNG.}
\hypertarget{native-639-lower}{}
\expandafter\def\csname HistoricalNoteText35\endcsname{Der Kollimationsfehler des Königsberger Meridiankreises beträgt beinahe 2 Grad. Gewöhnlich bringen die Instrumentenbauer Vorrichtungen an den Instrumenten an, um die Kollimationsfehler zu beseitigen. Der Astronom dankt es ihnen nicht. Dass sich das Fadennetz bei einigen Meridiankreisen auch in vertikaler Richtung bewegen lässt, sollte man nicht als Vorzug ansehen.\hfill G.}
\expandafter\def\csname HistoricalNoteText36\endcsname{Bei der ersten Beobachtung habe ich mich an den Stundenwinkel gehalten, da die Rektaszension offenbar durch einen Druck- oder Schreibfehler entstellt ist und statt $357^\circ\,49'\,14''5$ $357^\circ\,49'\,4''5$ lauten sollte.}
\begin{center}\small Astronomische Nachrichten. Band III. Nr. 53. Seite 77..84. 1824 März 3.\\[3pt]\rule{.72\linewidth}{.3pt}\end{center}
\begin{center}\emph{Ehrenrettung.}\end{center}

In einer der letzten Ausgaben der in Genua erscheinenden Correspondance Astronomique wird Professor \textsc{Pasquich} beschuldigt, diejenigen Positionen des Kometen von 1821, die in Nr. 2 der Astr. Nachr. abgedruckt sind, nicht aus Beobachtungen am Äquatoreal abgeleitet zu haben, wie \textsc{Pasquich} selbst versichert hat, sondern sie erfunden, das heißt aus Elementen berechnet zu haben. Zur Begründung dieser Anklage wird nicht etwa behauptet, dass an den Tagen, für die die Kometenpositionen bekannt gemacht worden waren, an der Ofener Sternwarte gar nicht oder falsch beobachtet worden sei. Im Gegenteil, der Ankläger hat selbst an den Beobachtungen teilgenommen und legt selbst einen Teil der Originalbeobachtungen aus \textsc{Pasquich}s eigener, vom Notar beglaubigter Handschrift vor. Der Ankläger beweist auch nicht, und genau genommen behauptet er nicht einmal, dass die Originalbeobachtungen bei richtiger und vollständiger Reduktion andere Positionen ergäben als die von \textsc{Pasquich} bekannt gemachten. Worauf gründet sich also, wird man fragen, eine Anklage, mit der ein Untergebener seinen Vorgesetzten des entehrendsten Verbrechens bezichtigt, das man einem Astronomen vorwerfen kann? Lediglich darauf, dass das Instrument gar nicht justiert gewesen sei, dass aus Beobachtungen mit einem nicht justierten Instrument keine richtigen Positionen abgeleitet werden könnten und dass die von \textsc{Pasquich} bekannt gemachten Positionen, deren Richtigkeit durch die Rechnung von Dr. \textsc{Ursin} sofort bestätigt worden war, also auf unredliche Weise erfunden worden sein müssten.

Astronomen brauchen nur diese einfache Darstellung oder die Anklage selbst einigermaßen aufmerksam zu lesen, um zu erkennen, dass sie jeder Grundlage entbehrt. Jeder Astronom weiß, dass aus Beobachtungen mit einem nicht justierten Instrument ebenso zuverlässige Positionen abgeleitet werden können wie aus Beobachtungen mit einem vollkommen justierten. Man bestimmt durch geeignete Beobachtungen, wie stark das Instrument vom vollkommen justierten Zustand abweicht, und berücksichtigt die daraus folgenden Korrekturen bei allen anderen Beobachtungen. Die Rechnung ist stets die genaueste Art der Korrektur. Ob die \SourceLeaf{640} Korrekturen einige Sekunden oder einige Minuten betragen, ist für die Genauigkeit der Ergebnisse gleichgültig. Im Allgemeinen zieht man es aus Bequemlichkeit vor, es nur mit kleinen Korrekturen zu tun, und nicht wegen der Genauigkeit. Bei Abweichungen, die sich immer gleich stark auswirken, ist selbst das so gut wie belanglos\HistoricalNote{35}.

Die Behauptung des Anklägers, man könne aus Beobachtungen mit einem nicht justierten Instrument keine richtigen Positionen ableiten, läuft demnach nur auf ein Eingeständnis hinaus, das keines Kommentars bedarf. Seinem Vorgesetzten, der seit Langem als Mathematiker einen guten Ruf genießt, konnte natürlich eine Korrekturrechnung nicht schwerfallen, die ja schon einem einigermaßen geübten Anfänger möglich ist.

Wenn man die Anklage nur oberflächlich liest, könnte man vielleicht glauben, dass das Instrument sehr weit vom vollkommen justierten Zustand entfernt gewesen sei. Doch

\emph{Erstens} folgt dies nicht aus den angeblichen Beweisen. Es wird berichtet, dass \textsc{Pasquich} einmal in der Dämmerung den Kometen mit dem Äquatoreal nicht finden konnte, als der Ankläger ihn schon im Kometensucher sah. Das beweist gar nichts. Der Komet konnte tatsächlich im Gesichtsfeld sein und doch gerade wegen der Dämmerung bei der starken Vergrößerung noch nicht erkennbar sein, als ein lichtstarker Kometensucher ihn bei schwächerer Vergrößerung schon zeigte. Ebenso beweisen die Rechnungen des Anklägers zu den Beobachtungen für sich genommen noch gar nichts für eine mangelhafte Justierung. Tatsächlich hätte er mit \emph{solchen} Rechnungen eine Abweichung finden müssen, selbst wenn gar keine vorhanden gewesen wäre. Er vernachlässigt nämlich die Refraktion ganz und macht sich auf seine Art über \textsc{Pasquich}s Versicherung lustig, er habe sie bei der Reduktion der Beobachtungen berücksichtigt. Nach Behauptung des Anklägers könne die Refraktion im Maximum weder bei der Differenz der Rektaszensionen noch bei der der Deklinationen zwei Sekunden erreichen. Wusste der Ankläger auch diese leichte Rechnung nicht auszuführen? Die Differenz beträgt am 27. Februar bei beiden über eine Minute.

\emph{Zweitens} aber lässt sich aus den angeführten Beobachtungen tatsächlich hinreichend deutlich erkennen, dass die Abweichung des Instruments keineswegs so enorm groß war.

Die wenigen Beobachtungen von $\gamma$ Pegasi, die der Ankläger uns zur Kenntnis gebracht hat, reichen zwar für eine genaue und vollständige Bestimmung der Korrekturelemente des Instruments nicht aus; sie genügen aber, um zu beweisen, dass die Abweichungen nicht so groß sind, dass sie das Auffinden von Sternen erschweren; sie genügen ferner, um die Kometenbeobachtungen selbst sehr genau zu reduzieren. Umso lieber teile ich die Ergebnisse meiner Berechnung hierzu mit, da daraus mit aller nur wünschenswerten Deutlichkeit hervorgeht, dass die Beschuldigung mehr als grundlos, dass sie falsch ist.

Aus der Auswertung aller 5 Beobachtungen von $\gamma$ Pegasi erhielt ich
\begin{center}\footnotesize\setlength{\tabcolsep}{3pt}
\begin{tabular}{@{}l r@{}}
Abstand des Instrumentenpols vom wahren Himmelspol & $=5'\;5''56$\\
Stundenwinkel des erstgenannten Pols . . . . . . . . . . . & $212^\circ\;27'\;27''$\\
Korrektur des Index für den Stundenwinkel . . . . . & $3'\;41''57$\\
für die Deklination . . . . . . . . . . . . . . & $4\;57.55$
\end{tabular}\end{center}

Ob und um wie viel die optische Achse von der Parallelität zum Deklinationskreis und dieser von der Rechtwinkligkeit zum Äquatorkreis abweicht, lässt sich aus den Beobachtungen eines einzigen Sterns nicht \SourceLeaf{641} bestimmen; auf die Reduktion der Kometenpositionen kann dies aber bei dem geringen Unterschied zwischen den Deklinationen keinen merklichen Einfluss haben.

Die 5 Beobachtungen von $\gamma$ Pegasi werden mit diesen Elementen folgendermaßen
dargestellt:

\begin{center}\footnotesize\setlength{\tabcolsep}{4pt}
\begin{tabular}{r|r|r|l|r}
 & \multicolumn{2}{c|}{Korr. d. St. W. wegen} & Korrig. beob. & Untersch.\\
 & des Instr. & der Refr. & Stundenwinkel & v. wahren \\
\hline
Febr. 20 & $-162''99$ & $+147''26$ & $81^\circ\,36'\,45''01$ & $+5''42$ \\
22 & $-174.81$ & $+95.54$ & $69\ \ 38\ \ 14.47$ & $+3.23$ \\
22 & $-168.76$ & $+116.59$ & $75\ \ 29\ \ 8.57$ & $+0.33$ \\
27 & $-174.97$ & $+95.56$ & $69\ \ 28\ \ 41.31$ & $+1.42$ \\
27 & $-168.59$ & $+117.84$ & $75\ \ 39\ \ 9.99$ & $-10.40$ \\
\end{tabular}
\end{center}

\begin{center}\footnotesize\setlength{\tabcolsep}{4pt}
\begin{tabular}{r|r|r|l|r}
 & \multicolumn{2}{c|}{Korr. d. Dekl. wegen} & Korrig. beobach- & Untersch.\\
 & des Instr. & der Refr. & tete Dekl. & v. d. wahren \\
\hline
Febr. 20 & $-97''54$ & $-147''48$ & $14^\circ\,11'\,19''98$ & $-0''63$ \\
22 & $-54.04$ & $-96.10$ & $14\ \ 11\ \ 21.86$ & $+1.42$ \\
22 & $-74.06$ & $-116.33$ & $14\ \ 11\ \ 17.61$ & $-2.83$ \\
27 & $-53.50$ & $-96.15$ & $14\ \ 11\ \ 20.35$ & $+0.33$ \\
27 & $-74.68$ & $-117.59$ & $14\ \ 11\ \ 21.73$ & $+1.71$ \\
\end{tabular}
\end{center}

Wendet man dieselben Elemente zur Reduktion der Kometenbeobachtungen an, so erhält
man

\begin{center}\footnotesize\setlength{\tabcolsep}{4pt}
\emph{Für den Stundenwinkel}
\end{center}

\begin{center}\footnotesize\setlength{\tabcolsep}{4pt}
\begin{tabular}{r|r|r|l|r}
 & \multicolumn{2}{c|}{Korrektur wegen} & Korrigierter & \shortstack{Abw. d. Mitt.\\v. \textsc{Pasquich}s}\\
 & des Instr. & der Refr. & Stundenwinkel & Angaben \\
\hline
Febr. 22 & $-171''15$ & $+101''34$ & $71^\circ\,48'\,50''19$ & \raisebox{-5.5pt}[0pt][0pt]{$-4''34$} \\
22 & $-165.09$ & $+124.21$ & $77\ \ 39\ \ 19.12$ & \\
26 & $-164.37$ & $+134.96$ & $79\ \ 28\ \ 15.59$ & \raisebox{-5.5pt}[0pt][0pt]{$+1.37$} \\
26 & $-161.37$ & $+154.53$ & $82\ \ 44\ \ 53.16$ & \\
27 & $-161.34$ & $+157.80$ & $83\ \ \ 9\ \ 56.48$ & \raisebox{-5.5pt}[0pt][0pt]{$+4.95$} \\
27 & $-157.98$ & $+191.42$ & $87\ \ 10\ \ 33.44$ & \\
\end{tabular}
\end{center}

\begin{center}\footnotesize\setlength{\tabcolsep}{4pt}
\emph{Für die Deklination}
\end{center}

\begin{center}\footnotesize\setlength{\tabcolsep}{4pt}
\begin{tabular}{r|r|r|l|r}
 & \multicolumn{2}{c|}{Korrektur wegen} & Korrigierte & \shortstack{Abw. d. Mitt.\\v. \textsc{Pasquich}s}\\
 & des Instr. & der Refr. & Deklination & Angaben \\
\hline
Febr. 22 & $-61''15$ & $-100''78$ & $14^\circ\,38'\,10''07$ & \raisebox{-5.5pt}[0pt][0pt]{$+7''42$} \\
22 & $-82.05$ & $-123.18$ & $14\ \ 38\ \ 24.77$ & \\
26 & $-89.01$ & $-134.43$ & $14\ \ 23\ \ \ 9.56$ & \raisebox{-5.5pt}[0pt][0pt]{$+0.91$} \\
26 & $-102.09$ & $-154.65$ & $14\ \ 23\ \ 18.26$ & \\
27 & $-103.80$ & $-158.18$ & $14\ \ 18\ \ 44.02$ & \raisebox{-5.5pt}[0pt][0pt]{$-5.58$} \\
27 & $-120.75$ & $-193.94$ & $14\ \ 18\ \ 34.82$ & \\
\end{tabular}
\end{center}

Für jeden, dem astronomische Berechnungen nicht fremd sind, müssen diese
Abweichungen der klarste Beweis sein, dass \textsc{Pasquich}s Positionen
wirklich durch ordnungsgemäße Reduktionsrechnung aus den Beobachtungen
abgeleitet waren. Freilich enthält, wie vorhin gezeigt, die ganze Anklage
keinerlei Grund, dies zu bezweifeln. Freilich gibt es auch ganz abgesehen
von dieser Anklage keinerlei Grund zu einem solchen Zweifel, der schon
allein deshalb unsinnig wäre, weil es viel weniger Arbeit kostet, die
wirklichen Beobachtungen zu reduzieren, als die Kometenorte aus Elementen
zu berechnen. Doch diesmal liefert die Beobachtung selbst tatsächlich den
offenkundigsten Beweis für den echten Ursprung der von \textsc{Pasquich}
bekannt gemachten Kometenpositionen. Für die meisten Leser wird schon die
Geringfügigkeit der oben gefundenen Abweichungen ein solcher Beweis sein.
Ein viel \SourceLeaf{642}\enlargethispage{1pt} stärkerer liegt aber noch in ihrem \emph{regelmäßigen Verlauf}.
Unsere Reduktionselemente können sich erheblich von den wahren unterscheiden.
Das muss einen gewissen, wenn auch stets nur geringen Einfluss auf die
Kometenpositionen haben. Vergleicht man die Ergebnisse der Rechnung nach
zwei verschiedenen Systemen von Reduktionselementen, die auf ein und dieselbe
Reihe von Beobachtungen angewandt werden, so zeigen sich kleine Unterschiede,
die aber notwendigerweise einem regelmäßigen Verlauf folgen. Vergliche man
hingegen Positionen, die aus Elementen berechnet wären, mit solchen, die aus
wirklichen Beobachtungen abgeleitet sind, so müssten sich Unterschiede
ergeben, an denen im Absoluten die Fehler der Elemente und im Relativen die
unregelmäßigen Beobachtungsfehler (zusammen mit etwa absichtlich und willkürlich
vorgenommenen kleinen Änderungen, falls man solche annehmen wollte) deutlich
zu erkennen wären. Es wäre mehr als ein Wunder, wenn jemand, der so unvernünftig
wäre, Positionen nach Elementen zu erfinden, obwohl er gute, leicht zu
reduzierende Beobachtungen vor sich hat, bei einem so törichten Betrug eine
solche Fünferkombination aus der Lostrommel zöge, dass er haargenau dasselbe
träfe, was ihm die Reduktion seiner Beobachtungen ergeben hätte. Wie regelmäßig
die oben genannten Unterschiede verlaufen, sieht man übrigens am besten, wenn
man die Reduktionselemente nicht aus den Beobachtungen von $\gamma$ Pegasi,
sondern aus denen des Kometen selbst ableitet, indem man sie mit den bekannt
gemachten Positionen vergleicht, und zwar nur aus zwei Beobachtungen des
Kometen, um dann zu sehen, wie die dritte damit übereinstimmt. Auf diese
Weise findet man\HistoricalNote{36}:

\begin{center}\footnotesize\setlength{\tabcolsep}{3pt}
\begin{tabular}{@{}l r@{}}
Abstand des Pols des Instruments vom wahren Himmelspol & $2'\;32''41$\\
Stundenwinkel des ersteren . . . . . . . . . . & $170^\circ\;25\;17''$\\
Korrektur des Index für den Stundenwinkel . . . & $-3\;23.47$\\
Korrektur des Index für die Deklination . . . . . & $-1\;34.38$
\end{tabular}\end{center}

Die auf ähnliche Weise wie oben durchgeführte Reduktion der Beobachtungen
des Kometen ergibt hiernach

\begin{center}\footnotesize\setlength{\tabcolsep}{3pt}\begin{tabular}{r | r r | r | r r | r}
 & \multicolumn{2}{c|}{Stundenwinkel} & nach \textsc{Pasquich} & \multicolumn{2}{c|}{Deklination} & nach \textsc{Pasquich} \\
\hline
Febr. 22 & $74^\circ\,44'$ & $8''99$ & $74^\circ\,44'\,9''$ & $14^\circ\,38'$ & $9''82$ & $14^\circ\,38'\,10''$ \\
26 & $81$ & $6\ \ 32.98$ & $81\ \ 6\ \ 33$ & $14$ & $23\ \ 13.36$ & $14\ \ 23\ \ 13$ \\
27 & $85$ & $10\ \ 10.03$ & $85\ \ 10\ \ 10$ & $14$ & $18\ \ 44.81$ & $14\ \ 18\ \ 45$ \\
\end{tabular}\end{center}

also tatsächlich \emph{vollkommene} Übereinstimmung, da \textsc{Pasquich}
keine Bruchteile von Sekunden angesetzt hat.

Dieses System von Reduktionselementen gibt die Beobachtungen von $\gamma$ Pegasi
hinreichend genau wieder, wenn auch nicht ganz so genau wie das obige.
Die Unterschiede betragen

\begin{center}\footnotesize\setlength{\tabcolsep}{4pt}\begin{tabular}{r | r | r}
 & \shortstack{im Stunden-\\winkel} & \shortstack{in der\\Deklination} \\
\hline
Febr. 20 & $+3''51$ & $-0''39$ \\
22 & $+12.47$ & $-10.32$ \\
22 & $+4.08$ & $-9.82$ \\
27 & $+10.81$ & $-11.42$ \\
27 & $-6.80$ & $-5.10$ \\
\end{tabular}\end{center}

Die von \textsc{Pasquich} tatsächlich angewandten Reduktionselemente kennen
wir nicht. Sie können sich erheblich von den eben angeführten unterscheiden,
da es zu unsicher ist, aus den Korrekturen auf \SourceLeaf{643} die ihnen
zugrunde liegenden Elemente zurückzuschließen, zumal wir die Korrekturen
für das Instrument aus den bekannt gemachten Daten nicht getrennt von der
Refraktion ermitteln können und \textsc{Pasquich} die Refraktionen vermutlich
nach anderen Tafeln und vielleicht nach anderen Methoden berechnet hat. Auch
gehört die Frage, ob die von \textsc{Pasquich} angewandten Reduktionselemente
die genauestmöglichen waren, gar nicht zur Sache. Oder vielmehr, schon diese
Frage bloß aufzuwerfen, heißt, \textsc{Pasquich} von dem gegen ihn erhobenen
Vorwurf freizusprechen. Ich erkläre jedoch gern, dass ich keinerlei Grund sehe,
daran zu zweifeln, dass \textsc{Pasquich} seine Reduktionselemente wohlüberlegt
aus wahrscheinlich viel zahlreicheren und vielleicht ganz anderen Beobachtungen
abgeleitet hat, da aus der obigen Berechnung selbst vollkommen klar hervorgeht,
dass \textsc{Pasquich} seine Beobachtungen tatsächlich als absolute und nicht
als Differenzialbeobachtungen reduziert hat.

Ich habe die Sache bisher bloß in wissenschaftlicher Hinsicht betrachtet.
Wie könnte man es aber unterlassen, sie auch unter dem Gesichtspunkt der Ehre
und Redlichkeit zu betrachten und die gerechte Empörung, die der edle
\textsc{Olbers} so treffend ausgesprochen hat, ganz zu teilen!

\noindent Göttingen 1824 März 3. \hfill C.\,F.\ \textsc{Gauss}.

[Handschriftliche Bemerkung: \textsc{Daniel Kmeth}, geb. in Bries (Brezno Bamja) in Ungarn 1783 Jan. 15, starb in Kaschau 1825 Juni 20.]
